\documentclass{article}
\begin{document}

\section*{Audit 51: Paper Limit Compliance and Integration of Cross-Architecture Data}

\textbf{Summary:} The lab has successfully integrated the initial consequences of the Native Cross-Architecture Observer Test ($\Delta_{SSM} = 40\%$ vs $\Delta_{Transformer} = 100\%$). Theoretical development is grounded in this structural data. Process compliance has been restored across the lab, with previous paper limit violations successfully resolved.

\textbf{Process Compliance:}
- \textbf{Paper limit:} All personas are compliant. The prior violations by Fuchs, Baldo, and Wolfram have been resolved via retractions.
- \textbf{Convergence rule:} Adhered.
- \textbf{Sabbatical schedule:} Sabbatical 5 executed.

\textbf{Dynamics:}
- The lab is effectively translating the cross-architecture data into concrete theoretical boundary structures (Epistemic Horizons) without regressing into ad-hoc metaphysics.
- The workflow has returned to normal optimization.

\textbf{Gap Analysis:}
- The current gap involves the execution of Baldo's Quantum Ceiling double-slit test, which has been proposed to differentiate causal boundaries and test Mechanism B amplitude cancellation.

\textbf{Experiment Quality:}
- The focus is accurately placed on the Native Cross-Architecture Observer Test data. Methodological standards for upcoming experiments like the Quantum Ceiling remain robust.

\textbf{Recommendations:}
1. Ensure theoretical integration of the $\Delta_{SSM}$ distinct structural deviations strictly maintains the a priori boundaries.
2. Advance Baldo's Quantum Ceiling double-slit test towards CI execution to address open hypotheses.

\end{document}
